\chapter{名词}
\section{三性四格二数}
\begin{itemize}
    \item 性 \\
          德语与英语区别之一在于，不仅人有性，物也有性。\textcolor{gray}{（英语物也有性的例外是：对船、车的称呼用she）}
          要明白性是德语名词不可分割的一部分
          \begin{myquote}
              背名词，必带冠。
          \end{myquote}
    \item 格
          \begin{itemize}
              \item 格的顺序
                    \begin{itemize}
                        \item 数字顺序：主属与宾，NGDA\textcolor{gray}{（语法书常用到第几格的说法）}
                        \item 课本顺序：主宾与属，NADG
                    \end{itemize}
              \item 名词变格 \\
                    名词的变格相对简单，黄书第一章就是它。
                    单数名词只有阳中的属格加-(e)s（多音节-s,单音节-es）。复数名词与格加-n。阳性名词中有部分“弱名词”，除单数主格外一律加-(e)n。
                    \begin{myquote}
                        复与-n，阳弱-(e)n，单阳中属-(e)s。
                    \end{myquote}
                    \textcolor{gray}{（复与n,这个在冠词、人称代词还会见到！所以说德语有规可循！）}
          \end{itemize}
    \item 性数规律
    \item 可视化：\\
          性色格形
          \textcolor{gray}{（用颜色形状表示性、格，数算一种特殊性）} \\
          \renewcommand{\arraystretch}{1.5}{
              \begin{tabular}{|M{2.5em}|M{3em}|M{2.5em}|M{3em}|}
                  \hline
                  格 & 形                 & 性 & 色               \\
                  \hline
                  N & \square[color=white, border=1] & M & \square[color=green]  \\
                  \hline
                  A & \diam[color=white, border=1]   & N & \square[color=green] \\
                  \hline
                  D & \olive[color=white, border=1]  & F & \square[color=purple] \\
                  \hline
                  G & \tri[color=white, border=1]    & P & \square[color=red]    \\
                  \hline
              \end{tabular}
          }

\end{itemize}